\conclusion
В данной работе исследованы ключевые аспекты контроля качества 3D-печати методом FDM, включая анализ дефектов, их причины и методы устранения. Автоматизация процесса выявления дефектов с использованием машинного обучения и компьютерного зрения стала приоритетной задачей. Сравнительный анализ традиционных алгоритмов обработки изображений и нейросетевых моделей позволил определить наиболее эффективный подход к решению задачи.

В рамках работы разработано мобильное Android-приложение, предназначенное для автоматического распознавания дефектов, возникающих в процессе 3D-печати. Основной функцией приложения является анализ загружаемых изображений деталей и выявление на них возможных отклонений от нормы. Для реализации этой задачи была обучена сверточная нейронная сеть, способная детектировать наиболее распространённые дефекты, такие как «спагетти», «паутина», деламинация, геометрические искажения и другие. В процессе работы выполнен полный цикл подготовки модели: сбор и аннотирование изображений, аугментация данных, обучение нейросети и тестирование её на реальных примерах.

Кроме разработки алгоритмов детекции, была создана серверная часть на базе Python и Flask, обеспечивающая обработку изображений, передачу данных между мобильным клиентом и моделью, а также визуализацию результатов. После загрузки фотографии деталь проходит анализ, а на изображении отмечаются области с выявленными дефектами. Полученные результаты передаются обратно в приложение, где пользователь может просмотреть аннотированное изображение и подробное описание обнаруженных проблем.

Для эффективной работы модели была использована платформа Roboflow для аннотирования и аугментации данных. Для создания интерфейса и взаимодействия с сервером была выбрана среда разработки Android Studio. В ходе тестирования проводился анализ точности предсказаний нейросети, а также выявление возможных ограничений и их устранение.

% Новые компетенции (добавлены комментариями, источник: Kompetentsii_IVTm-24.txt)
% Учебная практика, проектно-технологическая практика
% УК-1 – Способен осуществлять критический анализ проблемных ситуаций на основе системного подхода, вырабатывать стратегию действий.
% УК-3 – Способен организовывать и руководить работой команды, вырабатывая командную стратегию для достижения поставленной цели.
% ПК-4 – Способен проектировать сложные пользовательские интерфейсы.
% ПК-5 – Способен разрабатывать требования и проектировать программное обеспечение.
% ПК-6 – Способен разрабатывать пользовательские документы, а также стандартные технические документы на основе предоставленного материала.

При выполнении данной работы были освоены следующие компетенции:

УК-2 Способен определять круг задач в рамках поставленной цели и выбирать оптимальные способы их решения, исходя из действующих правовых норм, имеющихся ресурсов и ограничений. При выполнении работы были определены ключевые задачи и проанализированы способы их решения, проведено тщательное их сравнение и выбран наилучший подход.

УК-4 Способен осуществлять деловую коммуникацию в устной и письменной формах на государственном языке Российской Федерации и иностранном(ых) языке(ах). Во время выполнения работы осуществлялась коммуникация с научным руководителем и коллегами в устной и письменной форме, что соответствует компетенции.

УК-6 Способен управлять своим временем, выстраивать и реализовывать траекторию саморазвития на основе принципов образования в течение всей жизни. Разработка приложения потребовала планирования работы, распределения задач и изучения новых подходов в области машинного обучения и компьютерного зрения.

ОПК-1 Способен применять естественнонаучные и общеинженерные знания, методы математического анализа и моделирования, теоретического и экспериментального исследования в профессиональной деятельности. В процессе работы использованы принципы машинного обучения, компьютерного зрения и программирования для обработки изображений и распознавания дефектов.

ОПК-3 Способен решать стандартные задачи профессиональной деятельности на основе информационной и библиографической культуры с применением информационно-коммуникационных технологий и с учетом основных требований информационной безопасности. Во время выполнения работы были использованы различные информационно-коммуникационные технологии и ресурсы для поиска и анализа необходимой информации по теме работы.

ОПК-5 Способен инсталлировать программное и аппаратное обеспечение для информационных и автоматизированных систем. В ходе работы осуществлена настройка среды разработки, установка и конфигурация программного обеспечения для работы с нейросетью и API Roboflow.

ПК-1 Способен проводить анализ научной литературы, принимать участие в научно-исследовательских и опытно-конструкторских работах с использованием компьютерного моделирования, современных методов обработки информации и оформлять полученные результаты в виде отчетов, презентаций и докладов. В процессе выполнения работы проведён анализ научной литературы, применение полученных знаний на практике и использование современных методов обработки информации. Результат работы представлен в виде отчёта.

ПК-2 Способен производить анализ требований к программному обеспечению и проектировать программный продукт. В ходе работы сформулированы требования к функционалу приложения, спроектирована его архитектура и реализована серверная и клиентская часть для обработки изображений.

ПК-3 Способен создавать, модифицировать и сопровождать информационные системы. Разработано и протестировано приложение, интегрированное с нейросетью для детекции дефектов в 3D-печати, а также обеспечена его оптимизация для работы с различными устройствами.
